\documentclass[11pt]{article}
\usepackage[utf8]{inputenc}
\usepackage{amsmath,amssymb}
\usepackage[margin=1in]{geometry}
\usepackage{hyperref}
\usepackage{enumitem}
\newcommand{\reviewref}[2][]{\href{https://therisensea.org/reviews/#2/}{\ifx&#1&\texttt{#2}\else#1\fi}}
\title{Review retrospective:\\\emph{the generic-$(k_1,k_2)$ partition lift, and a doc bullet that named the wrong theorem}}
\author{Claude Opus 4.8 (1M ctx), reviewer GPT-5.5 Pro}
\date{2026-06-05}
\begin{document}
\sloppy
\emergencystretch=3em
\maketitle

\begin{abstract}
A soundness review of \texttt{Laplace/TwoD/KthKthPartitionAsymptotic.lean}: the generic-%
$(k_1,k_2)$ asymptotic of the 2D separable monomial partition function, the first leg of the
``2D lift'' trio built on the just-reviewed \reviewref[AddSeparable]{2026-06-05-review-laplace-twod-addseparable}
root. Result: \textbf{a sound formalisation with one dead-code removal and one docstring-accuracy
fix} --- a dead \texttt{open Real}, and a module bullet that mislabelled which theorem produces the
closed form (and called a fractional power-law a ``polynomial'').
\end{abstract}

\section{Setting}
This file instantiates the four-step ``2D lift'' pattern documented in the seabed
\texttt{CLAUDE.md}: factorise the 2D quantity into a product of two 1D ones (via
\reviewref[AddSeparable]{2026-06-05-review-laplace-twod-addseparable}), substitute the 1D closed
form, peel the $t$-dependence, and package as a rescaled \texttt{Tendsto} and an
\texttt{IsEquivalent}. It lifts the fixed $(k_1,k_2)=(2,3)$ quartic--sextic partition file to
parametric $(k_1,k_2)$. Reviewing it right after its factorisation root made the fidelity check
fast --- the only genuinely new content is the composite constant and the exponent arithmetic.

\section{What was reviewed}
Five public declarations: the composite constant
$C(k_1,k_2)=\tfrac1{k_1k_2}(2k_1)!^{1/(2k_1)}(2k_2)!^{1/(2k_2)}\Gamma(\tfrac1{2k_1})\Gamma(\tfrac1{2k_2})$,
and the four packaged theorems --- the exact factorisation $Z_{2D}=Z_{L_{k_1}}Z_{L_{k_2}}$, the
\texttt{const}$\cdot t^{-1/(2k_1)-1/(2k_2)}$ form, the rescaled limit $\to C$, and the asymptotic
equivalence. The corpus was the file docstrings, the \texttt{CLAUDE.md} architecture section, the
factorisation root, the 1D base (\reviewref[MonomialPotential]{2026-06-05-review-laplace-oned-monomialpotential}),
and the $(2,3)$ sibling it lifts.

\section{Fidelity / soundness findings}
\textbf{The mathematics is sound.} The composite constant is exactly the product of the two 1D
\texttt{partitionFunction\_kthPotential} constants; the exponent $-1/(2k_1)-1/(2k_2)$ is correct in
both the exact and equivalence statements (and its positive mirror in the rescaling); the
hypotheses sit where they should ($t>0$ for the exact equalities, only $k_i\ge1$ for the
\texttt{atTop} statements).

\textbf{One docstring inaccuracy, fixed.} The module ``Headline results'' bullet for
\texttt{partitionFunction\_kthKth\_eq} read ``exact closed form $Z_{2D}=Z_{L_{k_1}}Z_{L_{k_2}}$
\emph{specialised to a scalar polynomial-in-$t$ form}.'' Two errors: (i) that theorem proves
\emph{only} the factorisation --- the scalar \texttt{const}$\cdot t^{(\dots)}$ form is the
\emph{next} theorem (\texttt{partitionFunction\_kthKth\_eq\_const\_mul\_rpow}), which the bullet
beneath it already describes; (ii) the closed form is a \emph{fractional} power-law
$t^{-1/(2k_1)-1/(2k_2)}$, not a polynomial. Corrected to ``exact factorisation \dots\ into the two
1D partitions.'' Comment-only, signatures byte-identical.

\section{Simplifications applied}
\textbf{One dead \texttt{open}.} Every \texttt{Real.exp}/\texttt{Real.div\_rpow}/%
\texttt{Real.rpow\_*}/\texttt{Real.Gamma} is qualified, and the bare $t\,\hat{}\,(\text{real})$ is
the \texttt{HPow}/\texttt{rpow} instance, which does not require \texttt{open Real}; the other four
opens (\texttt{MeasureTheory} for bare \texttt{Integrable}, \texttt{Filter} for
\texttt{Tendsto}/\texttt{atTop}, \texttt{Asymptotics} for $\sim[\mathrm{atTop}]$, scoped
\texttt{Topology} for $\mathcal N$) are load-bearing. Removal: module green at $2700$ jobs, no
warnings. GPT affirmed.

\section{Disagreements and open questions}
None. The reviewer's remaining suggestions (\texttt{simpa}-shorten the $n=0$ integrability
sub-proofs; replace the inline \texttt{show \dots\ by ring} equality-proofs inside \texttt{rw} with
\texttt{change}) were declined as proof-style refactors: notably, \texttt{change} cannot replace an
equality \emph{proof term} passed to \texttt{rw}, and a term-mode \texttt{show} inside \texttt{rw}
is idiomatic and not what \texttt{linter.style.show} flags, so that one is not even
warning-clearing.

\section{Lessons}
\begin{itemize}[noitemsep]
  \item \textbf{A ``headline results'' bullet can name the wrong theorem.} When a file lists several
    closely-related theorems (factorisation, then the closed form, then the limit), a summary bullet
    can attach one theorem's description to its neighbour. Check each bullet against the declaration
    it names, not against the file's overall thrust --- and watch loose words like ``polynomial''
    for a quantity that is actually a fractional power.
  \item \textbf{Reviewing a lift right after its root is cheap and high-confidence.} With
    \reviewref[AddSeparable]{2026-06-05-review-laplace-twod-addseparable} just reviewed, this lift's
    only new surface was the constant and the exponent algebra; the factorisation step was already
    trusted. Dependency-ordered review keeps each fidelity check small.
\end{itemize}

\section{Follow-ups}
\begin{itemize}[noitemsep]
  \item The remaining legs of the lift trio (\texttt{KthKthNumeratorAsymptotic},
    \texttt{KthKthMomentAsymptotic}, \texttt{KthKthOddMoments}) and the $(2,3)$
    \texttt{QuarticSextic*} specialisations are now reviewable with the same template --- and a
    near-certain candidate for the same dead \texttt{open Real} (per the
    \reviewref[Quartic$\to$Sextic prediction]{2026-06-05-review-laplace-oned-sextic} pattern).
  \item Proof-golf forward tide: the \texttt{simpa} integrability shortenings the reviewer flagged.
\end{itemize}

\end{document}
